% \iffalse meta-comment
%
% Copyright (C) 2017- by Yanshuo Chu <yanshuoc@gmail.com>
%
% This file may be distributed and/or modified under the
% conditions of the LaTeX Project Public License, either version 1.3a
% of this license or (at your option) any later version.
% The latest version of this license is in:
%
% http://www.latex-project.org/lppl.txt
%
% and version 1.3a or later is part of all distributions of LaTeX
% version 2004/10/01 or later.
%
% \fi
%
% ^^A ======================================================================
% ^^A 用户手册附录：写作与排版建议、软件推荐、初砚硕课题组招聘。
% ^^A ======================================================================
%
% \changes{v3.2a}{2026/09/27}{自 v3.2a 起，示例文档完全参考官方的书写范例组织，
% 因此将写作与排版建议、Linux 软件列表及初砚硕课题组招聘挪至用户手册附录。}
%
% \chapter{写作与排版建议}[Writing and typesetting tips]
% \label{sec:tips}
%
% 本章收录数字书写与论文排版建议。
%
% \section{关于数字}[Numbers]
%
% 按《关于出版物上数字用法的试行规定》（1987年1月1日国家语言文字工作委员会等7个单位公布），
% 除习惯用中文数字表示的以外，一般数字均用阿拉伯数字。
%
% （1）公历的世纪、年代、年、月、日和时刻一律用阿拉伯数字，如20世纪，80年代，4时3刻等。
% 年号要用四位数，如1989年，不能用89年。
%
% （2）记数与计算（含正负整数、分数、小数、百分比、约数等）一律用阿拉伯数字，如3/4，4.5\%，
% 10个月，500多种等。
%
% （3）一个数值的书写形式要照顾到上下文。不是出现在一组表示科学计量和具有统计意义数字中的
% 一位数可以用汉字，如一个人，六条意见。星期几一律用汉字，如星期六。邻近两个数字并列连用，
% 表示概数，应该用汉字数字，数字间不用顿号隔开，如三五天，七八十种，四十五六岁，一千七八百元等。
%
% （4）数字作为词素构成定型的词、词组、惯用语、缩略语等应当使用汉字。如二倍体，三叶虫，
% 第三世界，“七五”规划，相差十万八千里等。
%
% （5）5位以上的数字，尾数零多的，可改写为以万、亿为单位的数。一般情况下不得以十、百、千、
% 十万、百万、千万、十亿、百亿、千亿作为单位。如345\,000\,000公里可改写为3.45亿公里或
% 34\,500万公里，但不能写为3亿4\,500万公里或3亿4千5百万公里。
%
% （6）数字的书写不必每格一个数码，一般每两数码占一格，数字间分节不用分位号“，”，凡4位或
% 4位以上的数都从个位起每3位数空半个数码（1/4汉字）。“3\,000\,000”，不要写成“3,000,000”，
% 小数点后的数从小数点起向右按每三位一组分节。一个用阿拉伯数字书写的多位数不能从数字中间转行。
% 示例宏包里的 \cs{num} 照这一条分组，见第~\ref{sec:example-package}~节。
%
% （7）数量的增加或减少要注意下列用词的概念：1）增加为（或增加到）过去的二倍，即过去为一，
% 现在为二；2）增加（或增加了）二倍，即过去为一，现在为三；3）超额80\%，即定额为100，现在
% 为180；4）降低到80\%，即过去为100，现在为80；5）降低（或降低了）80\%，即原来为100，现在
% 为20；6）为原数的1/4，即原数为4，现在为1，或原数为1，现在为0.25。应特别注意在表达数字减小
% 时，不宜用倍数，而应采用分数。如减少为原来的1/2，1/3等。
%
% \section{作图}[Drawing]
%
% \subsection{Ti\emph{k}Z}[TikZ]
% 使用 \pkg{tikz} 作图的核心思想是把格式、主题、样式与内容分离，定义在全局中。字少用五号字，
% 字多用小五。用 \pkg{tikz} 作图不会出现字体问题，图里的字体自动与正文一致。示例宏包里的
% \cs{rcell}、\cs{tcellone} 等几个宏就是这样画染色体示意图的，可以当例子看。
%
% \pkg{tikz} 适合画不需要大量实验数据支撑的示意图。对排版有极致追求的，可以参考
% \url{http://www.texample.net/tikz/resources/} 中所列工具，几乎所有作图软件作的图都可以转成
% \pkg{tikz}，然后在 \pkg{tikz} 中修改图中内容、定义字体，做到规范要求的图中字体一致。
%
% \subsection{R 作图}[Drawing with R]
% R 是一种极具代表性的作图工具，应用广泛。R 作图分两种情况：一是可以转换为 \pkg{tikz} 码，
% 二是不可转换。第一种情况图形简单，图中不含很多数据点，使用 R 语言中的 \pkg{tikz} 包即可。
% 第二种情况图形复杂，含有海量数据点，这时候不要转成 \pkg{tikz} 矢量图，那会使论文体积巨大，
% 推荐用 PDF 或 PNG 位图。用位图时要选好字号（五号或小五）和字体（宋体、新罗马），然后选定
% 生成的图的大小；在正文中用 \cs{includegraphics} 导入时，不要像导入矢量图那样控制图的大小，
% 用图原本的宽度和高度，这样位图里的文字才与正文一致。
%
% \section{源码与算法}[Source code and algorithms]
% 源码与算法都不在规范要求中，却常常用到（至少计算机学院），而且格式乱，每个实验室的要求都不一样。
% 模板照 iota-hit 把两者都当作一种图：文档类装 \pkg{listings} 与 \pkg{algorithm2e}，编号、题注、
% 外框、行号做成预设，见第~\ref{sec:listing}、\ref{sec:algorithm}~节，示例附录里各有例子。
% 要语法高亮更好的效果可以用 \pkg{minted}，但编译时要加 \texttt{-shell-escape} 选项并安装
% \texttt{pygmentize}，版式由文档自己管。
%
% \section{右翻页}[Open right]
% 对于双面打印的论文，可以让每章的标题页出现在右手边。规范中没有明确规定是否右翻页，
% 模板按部件给了开关，见 \option{layout} 族的 \option{openright}（第~\ref{sec:tristate}~节）。
%
% \section{术语与字体}[Terms and fonts]
% 缩略语交给 \cs{abbr}：首次出现全写、之后缩写，缩略语表与索引自动登记，见
% 第~\ref{sec:abbr}~节。思源字体（Noto CJK）由字体方案给出，见第~\ref{sec:chinese-fonts}~节。
%
% \section{编辑器与排版原则}[Editors and typesetting rules]
% \hologo{TeX} 源码编辑器推荐：（1）付费软件 WinEdt；（2）免费软件 Kile；（3）vim、Emacs、
% Sublime 等需要配置的编辑器。
%
% 格式和内容分离是 \hologo{LaTeX} 最大的优势，所有多次出现的内容、样式都可以定义为简单的
% 命令、环境，方便修改和管理。例如想把所有表示向量的符号由粗体 \cs{mathbf} 换成花体
% \cs{mathcal}，只需修改该格式命令的定义，不需要像 Word 那样处处修改。使用自定义命令和环境
% 才是正确使用 \hologo{LaTeX} 的方式。
%
% \chapter{这个星球上最好的免费 Linux 软件列表}[List of the best Linux software in our planet]
%
% \section{系统}[System]
%
% \href{http://fvwm.org/}{FVWM}，自从上世纪诞生以来，此星球最强大的窗口管理器。
% 推荐基于 FVWM 的桌面设计 hifvwm：\url{https://github.com/dustincys/hifvwm}。
%
% \subsection{hifvwm 的优点}[Advantages of hifvwm]
%
% \begin{enumerate}
%   \item 即使打开上百个窗口也不会“蒙圈”，对比 Windows 或 macOS 都无法做到。计算机性能越来越
%     强大，窗口任务的管理必须要升级到打怪兽级别。
%   \item 二维可视化任务栏。
%   \item 自动同步 Bing 搜索主页的壁纸。每次电脑开机、午夜零点自动更新，用户也可以手动更新，
%     从此审美再也不疲劳。
%   \item 切换窗口自动聚焦到最上面的窗口。用键盘快捷键切换窗口时减少操作过程，自动聚焦到
%     目标窗口。这一特性是虚拟窗口必须的人性化设计。
%   \item 类似 Windows 右下角那个显示桌面的功能：最小化窗口，在一个桌面之内操作多个任务。
%   \item 任务栏结合标题栏，节省空间。
%   \item 同类窗口切换，可以在同类窗口之内用类似 alt-tab 的方式切换。
%   \item ……
% \end{enumerate}
%
% \section{其他}[Others]
%
% \href{https://github.com/goldendict/goldendict}{GoldenDict}，星球最强大的桌面字典。
%
% \href{https://github.com/yarrick/iodine}{iodine}，“HIT-WLAN + 锐捷”时代的福音。
%
% \href{https://www.ledger-cli.org/}{ledger}，前“金融区块链”时代最好的复式记账系统。
%
% \href{https://orgmode.org/}{Org mode}，最强大的笔记系统，从来没有之一。
%
% \href{https://www.jianguoyun.com/}{坚果云}，国内一款支持 WebDAV 的云盘系统。
%
% \href{https://notmuchmail.org/}{notmuch}，目前最好的邮件管理工具，还在为每天几百个 email
% 苦恼？几百个这些都不算多，notmuch。
%
% \section{vim}[vim]
% 实现中英文每一句一行，以及每一句折叠断行的简单正则式，\hologo{TeX} 源码更加乖乖：
% \begin{latex}
% vnoremap <leader>fae J:s/[.!?]\zs\s\+/\="\r".matchstr(getline('.'), '^\s*')/g<CR>
% vnoremap <leader>fac J:s/[。！？]/\=submatch(0)."\n".matchstr(getline('.'), '^\s*')/g<CR>
% vnoremap <leader>fle :!fmt -80 -s<CR>
% \end{latex}
%
% \chapter{哈尔滨工业大学初砚硕课题组招（聘）全日制硕士、博士、助理研究员}[Yanshuo Chu lab is recruiting]
%
% \section{导师简介}[Principal Investigator (PI) introduction]
%
% 教授，博士生导师，国家高层次青年人才。2019年于哈尔滨工业大学获得计算机应用技术博士学位，
% 先后以博士后、数据科学家身份任职于美国MD安德森癌症研究中心基因医学系。2025年3月任职于哈尔滨
% 工业大学生命科学和医学学部（生物医学工程学科）和计算学部（智能科学与技术学科），两个学科都
% 招收硕博士。
%
% \section{主要研究方向}[Main research directions]
%
% 肿瘤信息学，肿瘤进化重构，肿瘤数字孪生。
%
% \section{学术成果}[Academic achievements]
%
% 近五年以第一作者在Nature Medicine（2023）、Nature Communications（2022）发表论文2篇，以
% 合作作者在Nature、Cancer Cell等业内顶级期刊发表论文10篇。其中Nature Medicine（2023）被
% Nature杂志特别报道，入选ESI高被引论文，荣获美国MD安德森癌症中心2023财年杰出科学奖。
%
% \section{团队介绍}[Research team introduction]
%
% 团队经费充沛、基础雄厚。所在团队于2019年申请并获批“生物大数据”教育部重点实验室，建立了
% 国内先进水平的基因组大数据专用高性能计算集群，拥有流式细胞仪、激光共聚焦显微镜等国际一流
% 水平的仪器设备，能够开展生物化学、分子生物学、细胞生物学、基因组学、生物信息学等各项实验。
%
% \section{其他介绍}[Other information]
%
% 强AI背景下，干实验室将走向无人化。目前我正基于哈尔滨医科大学医工交叉学院的计算集群开发和
% 部署自动化实验室管理系统，目标是开发首个无人化实验室。目前，该平台已能够超快速、大规模进行
% 单细胞、bulk、空间测序数据分析，极大缩短研究周期、快速得到研究成果。而在个人PC端，推荐
% 使用我开发的GPT2Org+elfeed工具链：
%
% \begin{itemize}
%   \item GPT2Org，一款
%     \href{https://chromewebstore.google.com/detail/gpt2org/nedljjoclmlgpobfohcnanjipnblgjbo?pli=1}{Chrome}、
%     \href{https://addons.mozilla.org/en-US/firefox/addon/gpt2org/}{Firefox}浏览器插件，使用
%     \href{https://platform.openai.com/api-keys}{OpenAI}或
%     \href{https://platform.deepseek.com/api_keys}{DeepSeek}的API，一键总结网页内容（学术
%     论文），并能够使用org-protocol自动发送到Emacs中，可以无缝整合org-roam知识库。
%   \item ElfeedAI，智能文献管理工具，可以自动更新论文并对论文重要性排序，进行AI阅读并总结。
%     我们不再需要去搜索文章，也不再需要详细阅读文章，这将极大节省注意力和时间，因为
%     “Attention Is All You Need”。
%     \href{https://www.bilibili.com/video/BV1dModYPE3c/?share_source=copy_web&vd_source=d6f2f3603fdf3b24c20399fbfd25df0d}{B站的演示视频《如何利用AI搞科研？》}
% \end{itemize}
%
% 以上，通过在计算集群端+个人PC端深度融合AI，将会有无限可能。
%
% \section{硕博士招生要求}[Graduate student recruitment requirements]
%
% \begin{itemize}
%   \item 计算机、数学、生命相关专业，熟悉常用编程语言。
%   \item 对科研和学术有一定的追求，有进取心，踏实肯干。
%   \item 发表过相关的高水平论文或有大规模软件、数学建模项目经验，优先。
% \end{itemize}
%
% \section{助理研究员招聘需求和待遇}[Assistant researcher position]
%
% 招聘数学、统计或计算机相关方向，工作地点在哈工大本部，招聘人数1-2人。负责：软件开发与算法
% 实现、参与科研项目的软件系统开发，完成算法设计、代码实现及性能优化；维护和升级现有科研工具
% 或仿真平台，确保其稳定性和扩展性；针对复杂问题建立数学模型，进行理论分析或数值模拟；跨学科
% 协作，与团队其他成员（如生命科学研究人员）合作，推动研究成果落地。
%
% 技能要求：熟练掌握Python/C++/MATLAB等至少一门语言，有软件开发经验；扎实的数学基础（如优化
% 理论、微分方程、概率统计等），具备数学建模经验；发表过SCI/SSCI论文者优先，或具备较强的
% 中英文技术文档撰写能力。
%
% 福利待遇：底薪25万/年+绩效。提供国内外学术会议交流机会与职业发展支持。
%
% \section{申请方式}[Application method]
%
% 发送邮件：\href{mailto:yanshuochu@hit.edu.cn}{yanshuochu@hit.edu.cn}
%
% \endinput
